\section{Receiver Configurations and Evaluation Protocol}
\label{app:evaluation-protocol}

\subsection{Learned Receiver Configurations}
\label{app:receiver-configurations}

The STMP baseline adapts \cite{cai2026} to the AFDM model, replacing
the Gaussian-message-based activity-probability update with learned
density-ratio estimation. It uses a white-noise-trained first-order score network, a
learned second-order score network, and a classifier for density ratio estimation.
GP-MP uses a denoiser trained with white Gaussian noise of covariance
$v\I_D$ and the scalar preconditioning
in Section~\ref{sec:network-settings}, its derivative-based variance
estimate, the corresponding activity classifier, and the input
variance update in Section~\ref{sec:variance-update}.
The two configurations use the same clean training channels and number of
denoiser training epochs.

GP-MP uses the white-noise-trained networks in
Figs.~\ref{fig:cstmp-denoising}--\ref{fig:cstmp-pilots}
and Table~\ref{tab:activity-evidence}; GP-MP-EP in
Figs.~\ref{fig:cstmp-snr} and \ref{fig:cstmp-pilots} uses the same networks.
The GP-MP results in Figs.~\ref{fig:cstmp-users}--\ref{fig:cstmp-cross} use earlier checkpoints
with different training perturbations and await replacement.

The activity-probability comparison in Section~\ref{sec:activity-evidence}
changes only the probability calculation within each receiver.
GP-MP-EP uses the same network weights and message-passing updates
as the corresponding known-phase GP-MP configuration, with phases
estimated once as described in Section~\ref{sec:phase-estimation}.

\subsection{Paired Scenes and Parameter Selection}
\label{app:paired-scenes}

We draw independent Bernoulli activity indicators with nominal probability
$p$ and retain scenes with at least one active device. The receiver uses the
resulting marginal probability $\rho=p/[1-(1-p)^K]$.

The evaluation protocol specifies $1{,}000$ Monte Carlo trials per
operating point. The counts below record the batches currently represented
in the working-draft figures and table; the remaining trials will be
completed before finalizing the corresponding results.

The comparisons in Section~\ref{sec:simulations} use saved, paired physical
scenes, with scene seeds reused across conditions.
The SNR figure uses scenes $0$--$999$ at each point for all
nine configurations. GP-MP-EP uses estimated phases, and the other eight
use known phases.
The pilot comparison reports channel NMSE for all nine configurations
at $Q\in\{8,12,16,20,24\}$, using $500$ paired scenes per point.
The device-count, activity-probability, and cross-profile
comparisons also use scenes $0$--$499$ at each condition.
The channel-denoising comparison uses $2{,}000$ channel samples from
$500$ scenes at each SNR. Both neural denoisers use the same channel
samples and noise realizations. The sparse-recovery baselines use an
identity observation operator and known AWGN variance, retaining their
receiver settings and GNN weights; GAMP-SBL retains its internal
noise-variance update.
The convergence comparison uses $30$ paired
scenes, and the activity-update comparison uses $100$ paired scenes per SNR.
Some evaluation scenes were used during method development and
selection of baseline activity thresholds; the adopted baseline activity decision rules are
reused without refitting. The reported SNR range was narrowed after
interim inspection; the earlier $8$ and $10$~dB results remain archived.
All finite outputs are included in the reported metrics, including those that
reach the iteration limit.

\subsection{Activity Probability Comparison}
\label{app:activity-evidence-protocol}

Table~\ref{tab:activity-evidence} currently uses the first $100$ completed
paired scenes per SNR.
All four variants use the same saved observations, known phases, and
noise variances estimated from the observations, with single-precision
inference and TF32 disabled.
The subset was adopted after interim inspection; the six SNR conditions
reuse the same physical scenes.

The limit rate is the fraction of runs that execute $600$ iterations.

\subsection{Device-Count-Specific Damping}

STMP uses damping factor $0.2$ at $K=20,40$ and $0.1$ at
$K=60,80,100$. The latter was selected after inspecting the original
$500$-scene results with damping factor $0.2$ and a diagnostic test with $12$ inputs at
damping factor $0.1$. All $12$ diagnostic inputs are included in the reported panels.
With the original damping factor, $498$, $479$, and $500$ runs reached the
iteration limit at $K=60,80,100$, respectively; those outputs remain
archived. Under the adopted settings, all $500$ runs satisfy the stopping criterion at
$K=20,40,60$, while $3$ and $14$ reach the limit at $K=80$ and $100$.

For GP-MP, the damping factor was reduced from $0.4$ to $0.2$ for all
$K=100$ scenes after preliminary inspection. At $K=80$, only scene~372
was rerun with damping factor $0.2$ after its original run reached the iteration
limit; this selected replacement is used in the reported panel. Other
$K=80$ scenes retain damping factor $0.4$. Both receivers retain the same
$60$--$600$ iteration limits and stopping criteria.
These damping factors were selected retrospectively.

\subsection{Convergence and Channel-Profile Evaluation}

Figure~\ref{fig:cstmp-convergence} uses the first $30$ paired scenes
at $0$~dB. For this experiment, GP-MP uses the AWGN-trained denoiser
and its matched activity classifier. GP-MP, STMP, and the baseline Off-grid SBL configuration
each execute $300$ iterations with early stopping disabled; the figure
shows the first $150$ iterations. At each iteration, NMSE is computed
from the total error energy divided by the total active-channel energy
over these scenes, followed by conversion to dB.

In the $1{,}000$-scene-per-SNR experiment displayed in
Fig.~\ref{fig:cstmp-snr}, STMP and GP-MP reach the iteration limit
in one and three runs, respectively, all at $-4$~dB. GP-MP-EP reaches
the limit in two runs, one at $-4$~dB and one at $-2$~dB. The remaining
outputs of these three configurations satisfy the stopping criterion.

The activity-probability and cross-profile sweeps retain damping factors $0.2$
for STMP and $0.4$ for GP-MP. All $2{,}500$ proposed-receiver
runs in the activity-probability experiment meet the stopping criterion. STMP reaches the limit in one
scene at each of $p=0.20$ and $0.25$.
The activity sampling and receiver prior are as defined in
Appendix~\ref{app:paired-scenes}; the horizontal axis denotes
the nominal probability $p$.

For the cross-profile comparison, all $12{,}000$ outputs of each learned
receiver satisfy the receiver stopping criterion.
The experiment varies the tabulated power delay profiles and
fading distributions at fixed geometry, pilots, and receiver settings,
without retraining or changing the stopping rules. Path gains remain fixed within each block, so the comparison concerns
the evaluated delay and fading profiles.

\subsection{Additional Implementation Settings}
\label{app:implementation-settings}

Each training or validation scene supplies four clean channel samples
before multiplication by $\alpha_k$. Denoiser model selection uses fixed
validation-noise realizations. The classifier is randomly initialized
and trained with the denoiser fixed.

For phase estimation, each L-BFGS-B refinement is restricted to a
differentiable interval of the waveform. Coordinate refinement starts
from zero powers and uses two sweeps in descending order of initial
pilot-correlation scores, with the current parameters retained among
the candidates.

The normalization constant is $s=0.2647966453$.
The tolerances for changes in activity probabilities and site parameters,
and for deviations from the moving average of activity probabilities
are $10^{-3}$, $0.1$, and $10^{-3}$, respectively, with
$\zeta_{\rm hist}=0.2$, $L_s=5$, and $L_d=10$.

The stopping rule monitors relative changes in $\widehat{\mathbf X}$ and $v$,
the maximum change in activity probability $\max_k|p_k^{(i)}-p_k^{(i-1)}|$, and
$d_{\rm var}^{(i)}=|v_{\rm tar}^{(i)}-v^{(i)}|/
\max\{v_{\rm tar}^{(i)},v^{(i)}\}$.
It also requires accepted site updates with stable precision
$(\tau_k^{\rm ext})^{-1}$ and natural parameter $\mathbf E_k/\tau_k^{\rm ext}$, together with a small difference from the previous exponential moving average, $\max_k|p_k^{(i)}-w_k^{(i-1)}|$.
After this comparison, update
$w_k^{(i)}=(1-\zeta_{\rm hist})w_k^{(i-1)}+\zeta_{\rm hist} p_k^{(i)}$, starting from
$w_k^{(0)}=\rho$.
After $I_{\min}$ iterations, all these conditions must hold for $L_s$
consecutive iterations and hard activity decisions must remain unchanged
for $L_d$ iterations; otherwise, return the last iterate at $I_{\max}$.

